%%=============================================================================
%% Requirementsanalyse
%%=============================================================================

\chapter{Requirementsanalyse}%
\label{ch:requirementsanalyse}

Aansluitend op het literatuuronderzoek, waarin de huidige stand van zaken binnen het probleemdomein werd onderzocht en enkele uitdagingen, opportuniteiten en bestaande oplossingen in kaart werden gebracht, wordt in dit hoofdstuk de requirementsanalyse besproken. Het doel van deze fase is het bepalen van de vereisten waaraan de \acrlong{poc} binnen dit onderzoek moet voldoen. De inzichten uit de literatuurstudie vormden hierbij het uitgangspunt voor het opstellen van een geprioriteerde lijst van functionele~(\ref{sec:fr}) en niet-functionele~(\ref{sec:nfr}) requirements.

\vspace{1em}

Om deze requirements verder te valideren en aan te vullen, werden interviews afgenomen met zowel experten als vertegenwoordigers van de doelgroep. De volgende personen namen deel aan deze interviews:

\begin{itemize}
    \item Co-promotor \textbf{dhr. G. De Bruyn} (KBC Bank \& Verzekeringen)
    \item Co-promotor \textbf{dhr. S. Hostyn} (Intigriti)
    \item \textbf{Mevr. J. De Smedt} (25-jarige junior developer)
    \item \textbf{Mevr. M. Adams} (21-jarige student, student laboratorium technieker)
    \item \textbf{Dhr. M. Dhondt} (22-jarige student)
    \item \textbf{Dhr. M. Maratta} (21-jarige student)
    \item \textbf{Dhr. S. Landtsheer} (21-jarige student)
\end{itemize}

De vragenlijst en de bijbehorende antwoorden van de geïnterviewden zijn opgenomen in bijlage~\ref{ch:interviewvragen}. In sectie~\ref{sec:interviews-inzichten} worden uitsluitend de belangrijkste bevindingen en daaruit voortvloeiende aanpassingen besproken. Op basis van deze resultaten wordt tot slot in sectie~\ref{sec:requirements} een definitieve, geprioriteerde lijst van functionele en niet-functionele requirements opgesteld.

\newpage

\section{Functionele requirements}
\label{sec:fr}

\vspace{1em}

In eerste instantie worden de functionele requirements (\acrshort{fr}'s) gedefinieerd. Deze vereisten beschrijven welke functionaliteiten de \acrshort{poc} kan bevatten en hoe deze zich dient te gedragen vanuit het perspectief van de eindgebruiker. Ze bepalen met andere woorden \textit{wat} de applicatie moet doen en welke interacties door de gebruiker moeten kunnen worden uitgevoerd. Om deze requirements op te stellen, wordt binnen dit onderzoek gebruikgemaakt van \textit{User Stories}.

\vspace{1em}

\subsection{User Stories}

\vspace{1em}

Vooraleer een geprioriteerde lijst van requirements opgesteld wordt, werden diverse user stories uitgewerkt. Deze beschrijven \textit{wie} een bepaalde functionaliteit wenst, \textit{wat} deze functionaliteit inhoudt en \textit{waarom} ze van belang is. Op basis van deze user stories wordt nadien een geprioriteerde lijst van \acrshort{fr}'s opgesteld aan de hand van de MoSCoW-prioriteringsmethode \autocite{AgileBusinessMoSCoW}.

\vspace{1em}

\subsubsection{Rol: Gebruiker (heeft iedereen)}

\vspace{1em}

\textit{Detectie}
\vspace{0.5em}
\hrule
\vspace{1.2em}

Als een \textbf{gebruiker}, \\
wil ik \textbf{vaak voorkomende dark patterns op streamingdiensten} (\ref{subsec:meest-voorkomend}) \textbf{in apps van streamingdiensten kunnen detecteren en bekijken}, \\
zodat \textit{ik weet wanneer ik tijdens het streamen mogelijks gemanipuleerd word}.

\vspace{1em}

Als een \textbf{gebruiker}, \\
wil ik \textbf{vaak voorkomende dark patterns op streamingdiensten} (\ref{subsec:meest-voorkomend}) \textbf{op websites van streamingdiensten kunnen detecteren en bekijken}, \\
zodat \textit{ik onder andere ook op aankoopschermen op de hoogte gebracht word van dark patterns}.

\vspace{1em}

Als een \textbf{gebruiker}, \\
wil ik \textbf{een korte uitleg kunnen raadplegen over een gedetecteerd dark pattern}, \\
zodat \textit{ik deze zelf beter kan herkennen}.

\vspace{1em}

Als een \textbf{gebruiker}, \\
wil ik \textbf{de mogelijke gevolgen van een gedetecteerd dark pattern kunnen bekijken}, \\
zodat \textit{ik begrijp waarom het als problematisch beschouwd kan worden en de impact ervan beter kan inschatten}.

\vspace{1em}

Als een \textbf{gebruiker}, \\
wil ik \textbf{kunnen zien welke businessdoelstelling achter een gedetecteerd dark pattern kan schuilgaan}, \\
zodat \textit{ik begrijp waarom deze techniek wordt gebruikt}.

\vspace{1em}

Als een \textbf{gebruiker}, \\
wil ik \textbf{geïnformeerd worden dat detecties indicatief zijn en niet 100\% accuraat}, \\
zodat \textit{ik de resultaten als ondersteuning gebruik in plaats van als absolute waarheid}.

\vspace{1em}

Als een \textbf{gebruiker}, \\
wil ik \textbf{een betrouwbaarheidsniveau zien bij een detectie}, \\
zodat \textit{ik kan inschatten hoe zeker de app is van de detectie}.

\vspace{1em}

Als een \textbf{gebruiker}, \\
wil ik \textbf{een ernstniveau (risico-indicator) bij detecties zien}, \\
zodat \textit{ik kan zien welke patronen het meest effectief en potentieel schadelijk zijn}.



\vspace{1.2em}

\textit{Notificaties}
\vspace{0.5em}
\hrule
\vspace{1.2em}

Als een \textbf{gebruiker}, \\
wil ik \textbf{tijdens interactie met een streamingdienst notificaties ontvangen van gedetecteerde dark patterns}, \\
zodat \textit{ik op het moment zelf geïnformeerd word over mogelijke manipulatieve ontwerpkeuzes}.

\vspace{1em}

Als een \textbf{gebruiker}, \\
wil ik \textbf{kunnen instellen wanneer ik notificaties ontvang}, \\
zodat \textit{ik niet gestoord word door te veel waarschuwingen}.


\vspace{1.2em}

\textit{Overzicht}
\vspace{0.5em}
\hrule
\vspace{1.2em}

Als een \textbf{gebruiker}, \\
wil ik \textbf{een overzicht van eerdere detecties kunnen raadplegen}, \\
zodat \textit{ik eerdere waarschuwingen later opnieuw kan bekijken}.

\vspace{1em}

Als een \textbf{gebruiker}, \\
wil ik \textbf{in de app kunnen zien hoeveel tijd ik besteed op streamingplatformen}, \\
zodat \textit{ik een mogelijk verslavend effect van dark patterns kan opmerken}.


\newpage

\textit{Instellingen}
\vspace{0.5em}
\hrule
\vspace{1.2em}

Als een \textbf{gebruiker}, \\
wil ik \textbf{kunnen kiezen welke dark patterns gedetecteerd worden}, \\
zodat \textit{ik het aantal waarschuwingen naar wens kan beperken}.

\vspace{1em}

Als een \textbf{gebruiker}, \\
wil ik \textbf{kunnen kiezen welke dark pattern categorieën gedetecteerd worden} (static, dynamic, in-between), \\
zodat \textit{de detectie afgestemd kan worden op de gewenste scope en betrouwbaarheid}.

\vspace{1em}

Als een \textbf{gebruiker}, \\
wil ik \textbf{toestemming kunnen geven voor de vereiste machtigingen van de app}, \\
zodat \textit{ik gebruik kan maken van functionaliteiten die de Accessibility Service API vereisen}.


\vspace{1.2em}

\textit{Overige}
\vspace{0.5em}
\hrule
\vspace{1.2em}

Als een \textbf{gebruiker}, \\
wil ik \textbf{zelf een mogelijk dark pattern kunnen melden}, \\
zodat \textit{mogelijk onopgemerkte patronen door de beheerders van de app onderzocht kunnen worden}.

\vspace{1em}

Als een \textbf{gebruiker}, \\
wil ik \textbf{een specifieke rol kunnen aanvragen}, \\
zodat \textit{ik kan helpen bij de correctheid en het succes van de app}.

\vspace{1em}

Als een \textbf{gebruiker}, \\
wil ik \textbf{juridische informatie en het privacybeleid kunnen raadplegen}, \\
zodat \textit{ik weet hoe mijn gegevens gebruikt worden}.



\vspace{1.5em}

\subsubsection{Rol: Handhaver}

\vspace{1em}

Als een \textbf{handhaver}, \\
wil ik \textbf{mij kunnen aanmelden}, \\
zodat \textit{ik gebruik kan maken van rolspecifieke functionaliteiten}.

\vspace{1em}

Als een \textbf{handhaver}, \\
wil ik \textbf{gedetecteerde dark patterns kunnen bekijken}, \\
zodat \textit{ik weet of een streamingdienst gebruikmaakt van dark patterns}.

\newpage

Als een \textbf{handhaver}, \\
wil ik \textbf{de bijhorende wetgeving bij een dark pattern kunnen raadplegen}, \\
zodat \textit{ik weet welke regelgeving van toepassing is}.

\vspace{1em}

Als een \textbf{handhaver}, \\
wil ik \textbf{kunnen zien welke businessdoelstelling achter een gedetecteerd dark pattern kan schuilgaan}, \\
zodat \textit{ik begrijp waarom deze techniek wordt gebruikt en de ernst van een overtreding kan inschatten}.

\vspace{1em}

Als een \textbf{handhaver}, \\
wil ik \textbf{een wijzigingsvoorstel kunnen doen van de bijhorende wetgeving bij een dark pattern}, \\
zodat \textit{verwijzingen naar wetgeving actueel blijven}.

\vspace{1.5em}

\subsubsection{Rol: Specialist}

\vspace{1em}

Als een \textbf{specialist}, \\
wil ik \textbf{mij kunnen aanmelden}, \\
zodat \textit{ik gebruik kan maken van rolspecifieke functionaliteiten}.

\vspace{1em}

Als een \textbf{specialist}, \\
wil ik \textbf{kunnen aangeven dat een detectie foutief is en correcte informatie kunnen meegeven}, \\
zodat \textit{de kwaliteit van de detecties verbeterd kan worden}.

\vspace{1em}

Als een \textbf{specialist}, \\
wil ik \textbf{aanvullende informatie kunnen toevoegen aan een detectie als voorstel}, \\
zodat \textit{gebruikers meer context krijgen}.

\vspace{1.5em}

\subsubsection{Rol: Beheerder}

\vspace{1em}

Als een \textbf{beheerder}, \\
wil ik \textbf{mij kunnen aanmelden}, \\
zodat \textit{ik gebruik kan maken van rolspecifieke functionaliteiten}.

\vspace{1em}

Als een \textbf{beheerder}, \\
wil ik \textbf{een overzicht kunnen raadplegen van foutieve detecties gemeld door specialisten}, \\
zodat \textit{ik de prestaties van het systeem kan evalueren en verbeteren}.

\newpage

Als een \textbf{beheerder}, \\
wil ik \textbf{voorstellen van aanvullende informatie over dark patterns kunnen beoordelen}, \\
zodat \textit{gebruikers meer context en een correcte uitleg ontvangen}.

\vspace{1em}

Als een \textbf{beheerder}, \\
wil ik \textbf{wijzigingsvoorstellen van juridische informatie door handhavers kunnen beheren}, \\
zodat \textit{verwijzingen naar wetgeving actueel en correct blijven}.

\vspace{1em}

Als een \textbf{beheerder}, \\
wil ik \textbf{rollen van gebruikers kunnen beheren}, \\
zodat \textit{ik specialisten en handhavers kan inschakelen om de app te verbeteren}.


\newpage

\subsection{Geprioriteerde Functionele Requirements}

\vspace{1em}

Legende: \enspace \textcolor{red}{\textbf{\bol}} Must have \enspace \textcolor{orange}{\textbf{\bol}} Should have \enspace \textcolor{ForestGreen}{\textbf{\bol}} Could have \enspace \textcolor{gray}{\textbf{\bol}} Won't have

\vspace{1em}

\textbf{Gebruiker+} \\ 
= Gebruiker met extra functionaliteiten (specialist, handhaver, beheerder)

\vspace{1em}

\begin{longtable}{%
        >{\raggedright\arraybackslash\small}p{2.5cm}
        >{\raggedright\arraybackslash\small}p{10.5cm}
        >{\centering\arraybackslash\small}p{1.5cm}
    }
    
    \caption[Geprioriteerde functionele requirements]{\label{tab:frs}Geprioriteerde lijst van functionele requirements op basis van de MoSCoW-methode} \\
   
    \toprule
    \textbf{Rol} & \textbf{Functionaliteit} & \textbf{Prioriteit} \\
    \midrule
    \endfirsthead
    
    \multicolumn{3}{c}%
    {\small{\tablename\ \thetable\ -- \textit{Vervolg van vorige pagina}}} \\
    \toprule
    \textbf{Rol} & \textbf{Functionaliteit} & \textbf{Prioriteit} \\
    \midrule
    \endhead
    
    \midrule \multicolumn{3}{r}{\small\textit{Vervolg op volgende pagina}} \\
    \endfoot
    
    \bottomrule
    \endlastfoot
    
    % =========================
    % OVERIGE
    % =========================
    
    Gebruiker &
    Aanvragen van een specialist- of handhaverrol &
    \textcolor{ForestGreen}{\bol} \\
    
    Gebruiker+ &
    Aanmelden &
    \textcolor{ForestGreen}{\bol} \\
    
    Gebruiker &
    Melden van mogelijk niet-gedetecteerde dark patterns &
    \textcolor{ForestGreen}{\bol} \\
    
    Gebruiker &
    Raadplegen van juridische informatie en privacybeleid &
    \textcolor{orange}{\bol} \\
    
    \multicolumn{3}{l}{} \\
    
    % =========================
    % DETECTIE
    % =========================
    
    \multicolumn{3}{l}{\small\textbf{Detectie}} \\
    
    Gebruiker &
    Detecteren van dark patterns in streamingapps (zowel in de mobiele app als website) &
    \textcolor{red}{\bol} \\
    
    Gebruiker &
    Raadplegen van een beschrijving van een gedetecteerd dark pattern &
    \textcolor{red}{\bol} \\
    
    Gebruiker &
    Raadplegen van de mogelijke gevolgen van een gedetecteerd dark pattern &
    \textcolor{orange}{\bol} \\
    
    Handhaver &
    Raadplegen van toepasselijke wetgeving bij dark patterns &
    \textcolor{ForestGreen}{\bol} \\
    
    Gebruiker &
    Raadplegen van de vermoedelijke businessdoelstelling achter een dark pattern &
    \textcolor{ForestGreen}{\bol} \\
    
    Gebruiker &
    Raadplegen van een betrouwbaarheidsniveau bij detecties &
    \textcolor{orange}{\bol} \\
    
    Gebruiker &
    Raadplegen van een ernstniveau bij detecties &
    \textcolor{ForestGreen}{\bol} \\
    
    Gebruiker &
    Disclaimer raadplegen over accuraatheid van de app &
    \textcolor{orange}{\bol} \\
    
    \multicolumn{3}{l}{} \\
    
    % =========================
    % NOTIFICATIES
    % =========================
    
    \multicolumn{3}{l}{\small\textbf{Notificaties}} \\
    
    Gebruiker &
    Ontvangen van notificaties wanneer een dark pattern wordt gedetecteerd &
    \textcolor{red}{\bol} \\
    
    Gebruiker &
    Configureren wanneer notificaties worden weergegeven &
    \textcolor{orange}{\bol} \\
    
    \multicolumn{3}{l}{} \\
    
    % =========================
    % OVERZICHT
    % =========================
    
    \multicolumn{3}{l}{\small\textbf{Overzicht}} \\
    
    Gebruiker &
    Raadplegen van een historiek van eerdere detecties &
    \textcolor{orange}{\bol} \\
    
    Gebruiker &
    Raadplegen van de tijd besteed op streamingplatformen &
    \textcolor{gray}{\bol} \\
    
    \multicolumn{3}{l}{} \\
    
    % =========================
    % INSTELLINGEN
    % =========================
    
    \multicolumn{3}{l}{\small\textbf{Instellingen}} \\
    
    Gebruiker &
    Instellen welke dark patterns gedetecteerd worden &
    \textcolor{ForestGreen}{\bol} \\
    
    Gebruiker &
    Instellen welke dark-patterncategorieën gedetecteerd worden &
    \textcolor{ForestGreen}{\bol} \\
    
    Gebruiker &
    Beheren van vereiste machtigingen voor detectiefunctionaliteit &
    \textcolor{red}{\bol} \\
    
    \multicolumn{3}{l}{} \\
    
    
    % =========================
    % MELDEN
    % =========================
    
    \multicolumn{3}{l}{\small\textbf{Melden}} \\
    
    Handhaver &
    Indienen van wijzigingsvoorstellen voor juridische informatie &
    \textcolor{ForestGreen}{\bol} \\
    
    Specialist &
    Melden van foutieve detecties met gecorrigeerde informatie &
    \textcolor{ForestGreen}{\bol} \\
    
    Specialist &
    Indienen van aanvullende informatie bij detecties &
    \textcolor{ForestGreen}{\bol} \\
    
    Beheerder &
    Raadplegen van foutieve detecties gerapporteerd door specialisten &
    \textcolor{ForestGreen}{\bol} \\
    
    Beheerder &
    Beoordelen van voorstellen voor aanvullende informatie &
    \textcolor{ForestGreen}{\bol} \\
    
    Beheerder &
    Beheren van juridische wijzigingsvoorstellen &
    \textcolor{ForestGreen}{\bol} \\
    
    Beheerder &
    Beheren van gebruikersrollen &
    \textcolor{ForestGreen}{\bol} \\
    
    \multicolumn{3}{l}{} \\
    
    %\bottomrule
\end{longtable}

\newpage

\section{Niet-functionele requirements}
\label{sec:nfr}

\vspace{1em}

Na het bepalen van een eerste versie van de \acrshort{fr}'s gaat deze sectie dieper in op het definiëren van niet-functionele requirements (\acrshort{nfr}'s) voor de \acrshort{poc}. In tegenstelling tot functionele requirements, die beschrijven \textit{wat} het systeem moet doen, leggen niet-functionele requirements de nadruk op \textit{hoe} het systeem dit moet uitvoeren. Denk onder meer aan aspecten zoals prestaties, betrouwbaarheid en gebruiksvriendelijkheid. 

\vspace{1em}

De \acrshort{nfr}'s in dit onderdeel worden opgesteld aan de hand van het \textit{ISO/IEC 25010:2011} productkwaliteitsmodel \autocite{ISO25010_2011}, zoals aangereikt binnen de opleiding als referentiekader voor het uitwerken van niet-functionele vereisten.
Verder worden de \acrshort{nfr}'s geformuleerd volgens het SMART-principe.

\vspace{1em}

\subsubsection{Functionele geschiktheid}
\vspace{1em}

\hspace*{1.5em}\textbf{Functionele compleetheid}

\begin{itemize}
    \item Minimaal 4 van de 8 vaak voorkomende dark patterns op streamingdiensten (\ref{subsec:meest-voorkomend}) moeten gedetecteerd kunnen worden.
    \item De applicatie moet dark patterns kunnen detecteren op zowel streamingwebsites als streamingapps.
    \item Alle functionele requirements met prioriteit Must Have moeten geïmplementeerd zijn in de \acrshort{poc}.
\end{itemize}

\vspace{0.5em}

\hspace*{1.5em}\textbf{Functionele correctheid}

\begin{itemize}
    \item Minimaal 80\% van de detecteerbare dark patterns moet correct geclassificeerd worden volgens de gebruikte taxonomie.
    \item De applicatie mag niet meer dan 15\% vals-positieve detecties genereren.
    \item De applicatie mag niet meer dan 20\% vals-negatieve detecties genereren.
\end{itemize}

\newpage

\subsubsection{Prestatie-efficiëntie}
\vspace{1em}

\hspace*{1.5em}\textbf{Snelheid}

\begin{itemize}
    \item Een detectieresultaat moet binnen 5 seconden beschikbaar zijn nadat een nieuw scherm verschijnt.
    \item Een melding moet binnen 1 seconde weergegeven worden nadat een dark pattern werd gedetecteerd.
\end{itemize}

\vspace{0.5em}

\hspace*{1.5em}\textbf{Capaciteit}

\begin{itemize}
    \item Minimaal 5 verschillende detecties moeten gelijktijdig uitgevoerd kunnen worden \textit{(bij snel wijzigen van een scherm)}.
    \item Minimaal 20 voorgaande detecties moeten opgeslagen kunnen worden in de detectiegeschiedenis.
\end{itemize}

\vspace{1em}

\subsubsection{Uitwisselbaarheid}
\vspace{1em}

\hspace*{1.5em}\textbf{Beïnvloedbaar}

\begin{itemize}
    \item De applicatie mag de normale werking van streamingapps niet verstoren.
\end{itemize}

\vspace{1em}

\subsubsection{Bruikbaarheid}
\vspace{1em}

\hspace*{1.5em}\textbf{Leerbaarheid}

\begin{itemize}
    \item Minimaal 75\% van de testgebruikers moet zonder hulp een detectie kunnen raadplegen.
    \item Minimaal 75\% van de testgebruikers moet zonder hulp meldingsinstellingen kunnen aanpassen, ook als ze voor het eerst met de applicatie werken.
\end{itemize}

\vspace{0.5em}

\hspace*{1.5em}\textbf{Voorkomen gebruikersfouten}

\begin{itemize}
    \item De applicatie moet gebruikers waarschuwen wanneer de juiste machtigingen voor Accessibility Services niet gegeven zijn.
\end{itemize}

\vspace{0.5em}

\hspace*{1.5em}\textbf{Volmaaktheid gebruikersinteractie}

\begin{itemize}
    \item Waarschuwingen mogen maximaal 25 woorden bevatten zodat de cognitieve belasting voor de gebruiker beperkt blijft.
    \item Elke detectiemelding moet een visuele risico-indicator bevatten.
\end{itemize}

\vspace{1em}

\subsubsection{Beveiligbaarheid}
\vspace{1em}

\hspace*{1.5em}\textbf{Authenticiteit}

\begin{itemize}
    \item Beheerders moeten zich authenticeren voordat beheerfunctionaliteiten beschikbaar zijn.
    \item Specialisten moeten zich authenticeren voordat validatiefuncties beschikbaar zijn.
    \item Handhavers moeten zich authenticeren voordat wijzigingsvoorstellen voor juridische informatie ingediend kunnen worden.
\end{itemize}

\vspace{1em}

\subsubsection{Onderhoudbaarheid}
\vspace{1em}

\hspace*{1.5em}\textbf{Modulariteit}

\begin{itemize}
    \item Detectielogica, gebruikersinterface en gegevensopslag moeten in afzonderlijke modules geïmplementeerd worden.
\end{itemize}

\vspace{0.5em}

\hspace*{1.5em}\textbf{Wijzigbaarheid}

\begin{itemize}
    \item Nieuwe wetgeving en uitleg over een patroon moet toegevoegd kunnen worden zonder wijzigingen door te voeren aan de detectielogica.
\end{itemize}

\vspace{1em}

\subsubsection{Betrouwbaarheid}
\vspace{1em}

\hspace*{1.5em}\textbf{Volwassenheid}

\begin{itemize}
    \item Tijdens de gebruikersevaluatie mag de applicatie geen kritieke crash vertonen.
    \item Minstens 90\% van de detecties moet succesvol verwerkt worden.
\end{itemize}

\vspace{0.5em}

\hspace*{1.5em}\textbf{Foutbestendigheid}

\begin{itemize}
    \item Wanneer een detectie mislukt door een netwerk- of interne fout, moet de applicatie binnen 5 seconden terugkeren naar een operationele toestand zonder herstart.
\end{itemize}

\vspace{0.5em}

\hspace*{1.5em}\textbf{Herstelbaarheid}

\begin{itemize}
    \item Minstens 95\% van de opgeslagen detectiegeschiedenis moet behouden blijven na een onverwachte applicatiecrash.
\end{itemize}

\vspace{1em}

\subsubsection{Overdraagbaarheid}
\vspace{1em}

\hspace*{1.5em}\textbf{Installeerbaarheid}

\begin{itemize}
    \item De applicatie moet compatibel zijn met Android-toestellen.
    \item De gebruiker moet de initiële configuratie van de applicatie binnen 5 minuten kunnen voltooien.
\end{itemize}

\vspace{1em}

\section{Randvoorwaarden}

\vspace{1em}

Naast het opstellen van de \acrshort{fr}'s en de \acrshort{nfr}'s, zijn er nog enkele aspecten die niet aan bod kwamen, maar waar wel rekening mee dient gehouden te worden bij het uitwerken van de \acrshort{poc} binnen de scope van dit onderzoek.

\begin{itemize}
    \item De ontwikkelde oplossing moet gebruikmaken van gratis of low-cost technologieën.
    
    \item De detectiefunctionaliteit moet gebaseerd zijn op een bestaande open-source detectietool.
\end{itemize}

Ten slotte wordt de scope van dit onderzoek beperkt tot de populairste streamingplatformen in Vlaanderen. De \acrshort{poc} richt zich daarom op platformen zoals Netflix, Disney+ en Streamz. Meer bepaald zal binnen dit onderzoek getest de \acrshort{poc} getest worden op het streamingplatform Netflix.

\newpage

\section{Kerninzichten uit de interviews}
\label{sec:interviews-inzichten}

\vspace{1em}

Om de requirements verder te verfijnen, werden interviews afgenomen met vertegenwoordigers van de doelgroep en twee experts. De doelgroep bestond uit vier studenten binnen de opleiding Toegepaste Informatica en een junior developer, terwijl de experten bestonden uit een systeemarchitect en een UI/UX-designer. De interviewvragen en antwoorden zijn terug te vinden in bijlage~\ref{ch:interviewvragen}. Op basis van deze gesprekken worden de belangrijkste inzichten hieronder kort samengevat.

\vspace{1em}

\subsubsection{Vertegenwoordigers van de doelgroep}

\vspace{1em}

\begin{itemize}
    \item Streamingdiensten worden voornamelijk gebruikt om op een \textbf{snelle} en \textbf{eenvoudige} manier content te bekijken.
    \item \textbf{Netflix} is het meest gebruikte streamingplatform.
    \item Er wordt voornamelijk gestreamd via computers/laptops, maar ook via smartphones en smart-tv's.
    \item Sommige respondenten waren niet vertrouwd met de term \textit{dark patterns}. De meeste konden wel een situatie schetsen waarin ze beïnvloed werden door een streamingplatform. De \textbf{autoplay}-functie werd hierbij het vaakst vermeld.
    \item Een tool die dark patterns detecteert moet \textbf{uitleg} geven over het gedetecteerde patroon, \textbf{waar} het zich bevindt, \textbf{waarom} het problematisch is en hoe men ermee kan omgaan. Daarnaast moet \textbf{personalisatie van waarschuwingen} mogelijk zijn.
    \item Waarschuwingen moeten vooral verschijnen bij situaties met grote gevolgen (bv. financiële impact), maar \textbf{subtiel blijven} zodat de gebruikerservaring op het streamingplatform niet wordt verstoord.
    \item De bereidheid om de tool te gebruiken daalt wanneer er te veel meldingen zijn, deze te opvallend zijn of interactie vereisen.
\end{itemize}

\vspace{1em}

\subsubsection{Experts}

\begin{itemize}
    \item ...
\end{itemize}

\vspace{1em}

\subsubsection{Aanpassingen}

\vspace{1em}

Op basis van deze bevindingen worden de volgende aanpassingen doorgevoerd om de requirements verder te verfijnen en het ontwerp beter af te stemmen op de noden van de gebruikers:

\begin{itemize}
    \item De \acrshort{fr} ``Raadplegen van een visualisatie waarop het dark pattern zichtbaar is'' wordt toegevoegd met prioriteit ``Must have''.
    \item De \acrshort{fr} ``Instellen van de frequentie van notificaties'' wordt toegevoegd met prioriteit ``Could have''.
    \item De prioriteit van de \acrshort{fr} ``Raadplegen van een betrouwbaarheidsniveau bij detecties'' wordt gewijzigd naar ``Must have''. Daarnaast wordt deze \acrshort{fr} aangevuld met de opmerking ``op basis van probability-percentage''.
    \item De \acrshort{nfr} ``Een melding moet binnen 1 seconde weergegeven worden nadat een dark pattern werd gedetecteerd.'' wordt weggehaald.
    \item De \acrshort{nfr} ``Minstens 95\% van de opgeslagen detectiegeschiedenis moet behouden blijven na een onverwachte applicatiecrash.'' wordt weggehaald.
    \item De \acrshort{nfr} ``Meldingen over gedetecteerde dark patterns mogen de gebruiker niet onderbreken tijdens het bekijken van content, enkel in het begin, op het einde of buiten de content.'' wordt toegevoegd onder de categorie ``Bruikbaarheid'' met als subcategorie ``Voorkomen gebruikersfouten''.
\end{itemize}

\newpage

\section{Definitieve requirements}
\label{sec:requirements}

\vspace{1em}

\subsubsection{Functionele Requirements}

\vspace{1em}

\subsubsection{Niet-Functionele Requirements}

\vspace{1em}